\section{Expert-guided graph ensembles and their intervention limits}
\label{exec20:expert-scope}

A bank building a diagram from several discovery algorithms will often face
agreement about a relationship but disagreement about its direction. The
methods may also miss different relationships because they use different
scores, tests and structural restrictions. Asking a domain expert to resolve
selected disagreements is a plausible use of an assistant. The design question
is what the expert is being asked, how its answer changes the graph, and whether
the resulting graph supports the bank's intervention query.

Tench and Demeester propose a dynamic expert-guided ensemble that distinguishes
adjacency from orientation. It counts how many component graphs contain any
edge between a pair, considers pairs in descending frequency, and adds directed
edges subject to an acyclicity check. A lower adjacency threshold determines
which pairs are eligible. For eligible pairs without a strict majority, it asks
the expert whether a plausible causal relationship exists. A separate threshold
controls when direction votes suffice and when expert orientation is requested.
The procedure uses equal component weights and a finite query budget. This is
greedy graph construction; the term model averaging does not establish a
Bayesian posterior over causal models. \citep{exec20_expert_averaging}

\subsection{An expert's causal ordering is not a direct edge}
\label{exec20:expert-semantics}

The source deliberately permits an expert to answer about a direct or indirect
causal path. This is weaker than knowing whether two nodes are adjacent in the
full graph. In a stipulated offer--uptake--balance chain
$O\rightarrow U\rightarrow Y$, changing the offer may affect balance through
uptake even though $O\rightarrow Y$ is not a direct edge. An expert can answer
the path question correctly while accepting a spurious direct connection
proposed by a component method. The paper recognizes this distinction in its
simulated expert and compares it with an expert that knows direct edges.
\citep{exec20_expert_averaging,exec20_expert_averaging_reference}

The proposed bank interface should therefore carry the proposition being
elicited: adjacency, ancestry, orientation conditional on adjacency, or a
temporal exclusion. It should also permit insufficient knowledge. A reply about
plausible ancestry cannot silently become proof of a direct edge, and a forced
binary answer cannot represent every legitimate uncertainty state. The set of
variables matters: marginalizing uptake changes the observed-variable graph
even when the underlying mechanism is unchanged.

The method's output is a DAG. Its component comparison includes FCI, while its
noisy benchmarks can contain latent confounding. A directed output cannot by
itself retain every bidirected or unresolved endpoint in a latent graph. Before
an adapter is proposed, the conversion from each component's graph type must be
explicit, with unsupported endpoints preserved or refused. Acyclicity checks
alone do not establish compatibility with an input equivalence class or identify
an intervention in the presence of hidden common causes.
\citep{exec20_expert_averaging}

\subsection{Frequency, diversity and structural compatibility}
\label{exec20:expert-votes}

For one binary question, suppose eleven voters each have probability $3/4$ of
being correct. Under independent errors their majority is correct with
probability
\begin{equation}
 \sum_{k=6}^{11}\binom{11}{k}
       \left(\frac34\right)^k\left(\frac14\right)^{11-k}
 =\frac{1012581}{1048576}\simeq0.966.
 \label{exec20:expert-majority}
\end{equation}
If all eleven instead copy the same $3/4$-correct signal, majority correctness
remains $3/4$. These are two authored probability laws with the same individual
accuracy. They show why individual competence alone cannot justify improvement
toward certainty. The paper itself acknowledges that algorithms share data and
often share statistical principles. Its discussion motivates diversity, but
does not estimate a calibrated graph probability from vote frequency.
\citep{exec20_expert_averaging,exec20_expert_averaging_reference}

Even individually valid DAGs can produce incompatible edge majorities. Consider
the three graphs
\begin{equation}
 \begin{split}
 G_1&=\{A\rightarrow B,\ B\rightarrow C\},\\
 G_2&=\{B\rightarrow C,\ C\rightarrow A\},\\
 G_3&=\{C\rightarrow A,\ A\rightarrow B\}.
 \end{split}
 \label{exec20:expert-cycle}
\end{equation}
Each is acyclic, yet every edge of $A\rightarrow B\rightarrow C\rightarrow A$
has frequency $2/3$. A greedy acyclicity rule must discard or reconsider an
edge. Its result can depend on tie handling and earlier choices. Returning one
DAG makes the output structurally admissible under that rule, but it neither
preserves every component claim nor proves which disputed edge is false.
\citep{exec20_expert_averaging_reference}

\subsection{The reported evidence is informative and conditional}
\label{exec20:expert-evidence}

The main experiments use seven benchmark network structures with synthetic
categorical samples, clean and corrupted versions, four sample sizes from 100
to 100,000, and twenty column orderings. The principal simulated expert is
correct on 70\% of existence questions and 90\% of orientation questions.
Thresholds of zero for adjacency and 0.9 for orientation, with a budget of twice
the node count, were selected using separate randomly generated networks. The
paper reports broad gains in its graph scores, alongside a precision--recall
tradeoff and varying comparison denominators when algorithms fail or return
undefined precision. These are reported experiments; they are not newly
reproduced results in this proposal. \citep{exec20_expert_averaging}

The appendix tests a narrower real-model setting on five clean networks,
using \texttt{gpt-5-nano-}\allowbreak\texttt{2025-08-07} and excluding Pathfinder
because its variable names lack the needed semantics and Formed because many
edges represent diagnostic rather than causal relations. Expert responses are
cached across seeds and sample sizes. Both argument orders are queried, with
disagreement causing a fallback to component voting. The reported per-response
orientation and existence accuracies are not the accuracy of that combined
fallback policy. \citep{exec20_expert_averaging}

\begin{center}
\begin{tabular}{lrr}
\hline
Reported case & Orientation accuracy & Existence accuracy\\
\hline
Mean across five networks & 93.3\% & 66.4\%\\
Property & 83.9\% & 43.9\%\\
Alarm & 82.8\% & 52.2\%\\
\hline
\end{tabular}
\end{center}
Property's existence accuracy falls below the paper's better-than-random expert
assumption. A useful mean therefore cannot replace domain-specific evidence,
and accuracy on all possible pairs need not equal accuracy on the difficult
pairs selected adaptively by an ensemble. Cached responses also do not become
twenty independent expert realizations because the numerical algorithms are
rerun twenty times. The appendix's near reconstruction of the Asia network
raises a memorization concern; unsuccessful recall of another full network
cannot establish that its concepts or fragments were absent from training.
\citep{exec20_expert_averaging}

The source states that it supplies no formal graph-recovery guarantee and has
not empirically qualified continuous or mixed observations. Its cost discussion
includes component discovery, expert requests and structural-distance scoring.
The proposed comparison should likewise retain timeout, memory and invalid-
output denominators, distinguishing a conditional score on successful outputs
from all-request performance. The source's proposed imputation sensitivity
analysis does not supply the missing runtime results or a bank operating cost.

\subsection{Resolve the printed query contract before reuse}
\label{exec20:expert-budget}

The prose says that a fixed budget bounds queries and exhausted decisions fall
back to majority voting. The printed Algorithm~1 leaves two details unresolved.
For an existence query, the counter increment lies inside the branch reached
only when the expert accepts the connection. With three eligible minority
connections, budget one, and three rejections, a literal trace makes three
expert calls while the printed counter remains zero. Separately, when an
accepted adjacency has two legal directions, neither reaches the orientation
threshold, and no budget remains, the printed conditional chain has no final
branch adding the majority direction. An example with direction votes
$3/4$ and $1/4$, threshold $9/10$ and budget zero therefore leaves the edge
unassigned under the printed chain. \citep{exec20_expert_averaging_reference}

These are specification discrepancies, not verified defects in the released
implementation. Resolving them requires an immutable source version and its
tests and call logs. A proposed adapter should count every attempted external
call, distinguish logical questions from the two argument-order requests,
account for retries and invalid answers, and define the behavior at exhaustion.
Caching needs a key that binds the variable meanings, full request, model
snapshot and relevant settings. Otherwise a cost-saving cache can silently
reuse evidence for a different scientific question. The method and data lead
in the paper are useful acquisition routes; package names and reported license
labels are not sufficient permission or provenance for direct adoption.

\subsection{A selected graph still has to answer the intervention query}
\label{exec20:expert-identification}

For a separate exact example, consider two standardized observational summaries
$X,Y$ whose temporal ordering has not been established. Let
$\rho=3/5$. The following two zero-mean Gaussian structural models use
independent errors within each model:
\begin{equation}
 \begin{aligned}
 \mathcal M_1:\quad &X=\epsilon_X,\quad Y=\rho X+\epsilon_Y,\\
 &\operatorname{Var}(\epsilon_X)=1,\quad
 \operatorname{Var}(\epsilon_Y)=1-\rho^2,\\[2pt]
 \mathcal M_2:\quad &Y=\eta_Y,\quad X=\rho Y+\eta_X,\\
 &\operatorname{Var}(\eta_Y)=1,\quad
 \operatorname{Var}(\eta_X)=1-\rho^2.
 \end{aligned}
 \label{exec20:expert-worlds}
\end{equation}
Both imply covariance
$\left(\begin{smallmatrix}1&3/5\\3/5&1\end{smallmatrix}\right)$ and hence the
same observational distribution. Their effects on the mean of $Y$ of increasing
the intervention level of $X$ by one unit differ:
\begin{equation}
 \tau(\mathcal M_1)=\frac35,\qquad \tau(\mathcal M_2)=0.
 \label{exec20:expert-effects}
\end{equation}
More algorithms fitted to the same distribution do not identify the missing
direction. Truthful temporal restrictions, suitable experiments or independently
justified structural knowledge can eliminate a candidate; that extra information
must be stated. The example does not override a known banking information clock.
\citep{exec20_expert_averaging_reference}

Suppose a one-unit increase in $Y$ is provisionally valued at USD~100 and the
intervention costs USD~40. The two compatible margins are USD~20 and USD~$-40$.
Treating a 90\% graph vote as probability for $\mathcal M_1$ gives effect
$27/50$ and margin USD~14, but the vote has supplied no probabilistic calibration
that warrants that calculation as expected value. A decision under unresolved
model ambiguity must retain the compatible effects or adopt an explicitly
justified probability model. This is an authored planning illustration, not
evidence of model performance or bank returns. \citep{exec20_expert_averaging_reference}

The proposal should therefore retain expert-guided averaging as a candidate
for generating and reviewing diagrams. A public prototype comparison can use
known generated mechanisms, withheld domains, correlated component failures,
wrong expert assertions, budget exhaustion and incompatible graph endpoints.
It should compare static aggregation, selective expert queries, a numerical
method without expert input and a workflow that preserves unresolved graph
classes. The outcomes should include the correctness of the intervention
answer, supported refusal, retained ambiguity, all physical requests and total
review cost. Graph scores remain useful diagnostics, but qualification of an
intervention service requires its own evidence.


\section{When an end-to-end graph workflow still lacks calibrated effects}
\label{rel64:graph-confidence}

The preceding example leaves a practical question. Can an assistant quantify
its uncertainty about each arrow, draw many graphs and pass those graphs to a
flexible generative model? Such a workflow would be useful if its final
intervals retained the mechanisms relevant to the bank's decision. Each step,
however, changes the object whose uncertainty is being described. A probability
assigned to a sentence, a distribution over graphs and a sampling interval for
an intervention effect require different justifications.

Molero P{\'e}rez's February 2026 thesis develops a concrete instance of this
workflow. Pairwise language judgments, optionally augmented by retrieval,
produce proposed causal relations. The study compares self-reported
confidence, selected-answer token probabilities and a HallBayes refinement
based on weakened prompts. Its chosen procedure converts the refinement's
risk score into an edge probability, averages three runs, samples 100 graphs
and fits causal neural spline flows. This is a substantive implementation
proposal with reported biomedical comparisons. Its availability permits a
more precise assessment than either accepting the integration as calibrated
or dismissing it merely because language models propose the graphs. The
following assessment concerns the printed method and results; it does not
reproduce the upstream software or its model calls. \citep{rel64molero}

\subsection{A success ceiling gives an error floor}

Write the Bernoulli relative entropy, in nats, as
\[
 d(p,q)=p\log\frac pq+(1-p)\log\frac{1-p}{1-q},
 \qquad 0<q<1.
\]
The thesis's equations (14)--(16) average clipped positive log-probability
lifts for the selected answer under original and masked prompts. It then uses
this quantity as an information budget in a Bernoulli inequality and defines
the largest compatible success probability. For any stipulated nonnegative
budget $\Delta$, the logical implication of that inequality is
\begin{equation}
 \begin{split}
 d(p,q)\leq\Delta,\qquad
 p_{\max}(\Delta,q)&=\sup\{u\in[0,1]:d(u,q)\leq\Delta\},\\
 p\leq p_{\max}(\Delta,q)&\quad\Longrightarrow\quad
 1-p\geq1-p_{\max}(\Delta,q).
 \end{split}
 \label{rel64:error-floor}
\end{equation}
The complement is therefore a lower bound on error under this premise. The
thesis calls it an upper bound and subsequently uses it to set graph-sampling
probabilities. That interpretation does not follow from the displayed
inequality. With $q=.1$ and $\Delta=.5$, the upper endpoint is approximately
$.49505$, giving complement $.50495$. Yet $p=.1$ satisfies the inequality
and has error $.9$. This exact counterexample concerns the direction of the
bound; it does not assert that every estimated graph in the thesis is wrong.
\citep{rel64molero}

There is a further distinction between the score and the premise. Evaluating
a clipped log-lift at one selected answer does not generally evaluate an
expectation over the full answer distribution. Nor does masking a prompt
establish the probability that its causal assertion is true in the intended
bank population. A separate calibration argument is needed before the score
can stand for the information budget or the causal correctness probability
in~\eqref{rel64:error-floor}. Averaging repeated calls reduces some
Monte Carlo variation conditional on the response mechanism; it cannot by
itself supply either missing interpretation.

\subsection{The information inequality remains useful in its stated direction}

The thesis cites the 2025 version of Chlon and colleagues' information-budget
paper. Its Theorem 5 follows from convexity and reduction to an event. Let
$P$ be a reference distribution, $S_k$ predictive distributions, and
$\bar S=\mathbb E_k S_k$ their mixture. For an event $A$, put
$p=P(A)$ and $q=\bar S(A)$. Under absolute continuity and finite expected
relative entropy,
\begin{equation}
 \mathbb E_k\operatorname{KL}(P\Vert S_k)
 \geq\operatorname{KL}(P\Vert\bar S)
 \geq d(p,q).
 \label{rel64:necessary-information}
\end{equation}
The first step is Jensen's inequality in the second distribution; the second
is the information lost when the outcome is reduced to membership of $A$.
Equality for the entire chain requires equality in both steps. Sufficiency
of the event for distinguishing $P$ from the mixture addresses the second
step; it does not eliminate a positive mixture Jensen gap.
\citep{rel64chlonone,rel64chlonthree}

A necessary information requirement is not a success guarantee. Consider the
project's four-outcome construction
\[
 S=(.05,.05,.001,.899),\qquad
 P=(.05,.05,.899,.001),\qquad A=\{1,2\}.
\]
Use the same $S$ for every prompt. Both event probabilities are $.1$, but
\[
 \operatorname{KL}(P\Vert S)=.898\log899\simeq6.10755
 >d(.95,.1)\simeq1.99421.
\]
The extra divergence lies within the failure event. Exceeding the information
requirement associated with a target success probability of $.95$ has not
raised actual reference-event probability above $.1$. This is a failure of a
sufficiency interpretation, not a counterexample to
inequality~\eqref{rel64:necessary-information}.

The June 2026 version changes the title, author list and deployment
description. It explicitly requires a specified Bernoulli predicate and a
reference supplied by gold labels, a deterministic verification primitive or
an audited calibrated verifier. Unrestricted generation without such a
reference is outside its stated regime. Its reported held-out audit remains
one particular evaluation, not an intervention-coverage result. These
qualifications improve the interpretation of the proposal, but they do not
turn the necessary inequality into a sufficient guarantee. Even with a
binary reference, $d(.05,.5)=d(.95,.5)$: without a justified branch
restriction, inversion does not distinguish low from high event probability.
With constant predictive probability $.5$ and target $.95$, the printed
budget ratio is one for either reference probability, before clipping matters.
The event, reference and direction of the bound must therefore remain visible
when adapting the planner. A later version's verifier assumptions also cannot
be silently attributed to the thesis's earlier masking procedure.
\citep{rel64chlonthree}

\subsection{Independent edge draws need a coherent graph law}

For an accepted edge the thesis uses one minus the risk score; for a rejected
edge it uses the risk score. It then samples edges independently and applies
topological sorting. Sorting represents an already acyclic graph in an
appropriate order; it cannot make every independently drawn graph acyclic.
The treatment of a failed sort is part of the probability law. The following
example diagnoses the general construction and does not claim that the
thesis's particular draws contained a cycle. \citep{rel64molero}

Suppose the only possible arrows are $A\to B$, $B\to C$ and $C\to A$,
each assigned probability $4/5$. Under independent draws,
\[
 \Pr(\text{cycle})=(4/5)^3=64/125,\qquad
 \Pr(A\to B\mid\text{acyclic})=\frac{36}{61}\simeq.59016.
\]
The conditional numerator is $(4/5)[1-(4/5)^2]$ and the denominator
is $1-(4/5)^3$. Rejecting cyclic draws changes the edge probabilities;
deleting an edge instead creates another distribution determined by the
deletion rule. Indeed, no distribution supported on DAGs can retain all
three proposed marginals: a DAG contains at most two of these arrows, whereas
their proposed expected count is $12/5$. This strengthens the majority-cycle
example in Section~\ref{exec20:expert-votes}: structural validity and
marginal probability interpretation cannot be checked independently.

A usable alternative defines a joint law on an explicitly admissible model
class, records its prior and likelihood interpretation, and tests sensitivity
to that class. When those probabilities cannot be justified, retain compatible
graphs and propagate their effect sets instead. Neither route assumes that a
linguistic probability score is already a posterior over banking mechanisms.

\subsection{An interval comparison needs an intervention reference}

The thesis trains four groups of 100 flow models: an expert reference graph,
the sampled language-derived graphs, the expert graph's ordering without its
adjacency restrictions, and random orderings without those restrictions.
Within each group, the 2.5th and 97.5th percentiles of the fitted effect
estimates define an interval. The reported interval-overlap means are .723,
.712 and .301 for the last three groups compared with the fitted expert-graph
group. The first two differ by .011. That contrast describes these fitted
model intervals; the table does not provide a paired uncertainty interval for
the difference. It is not evidence that language assistance generally
outperforms knowledge of the causal ordering. \citep{rel64molero}

Agreement with a fitted expert reference is useful when the task is to
reconstruct that reference. It is not observed intervention truth. If two
methods return the identical interval $[.5,.7]$ while the true effect is zero,
their interval overlap is one and both miss the effect. More generally,
percentiles across graph choices and training runs describe variation under
that procedure. A 95\% repeated-sampling coverage claim needs a defined
sampling experiment, valid treatment of graph uncertainty and an identified
target. A calibrated Bayesian interpretation instead needs a justified joint
posterior and corresponding model assumptions. Repetition alone supplies
neither interpretation.

The numerical comparator also needs matched information and a suitable test.
The thesis specifies chi-square conditional-independence tests for its PC
baseline while prioritizing continuous clinical variables. Its language
workflow additionally draws on outside knowledge. A fair follow-up would
resolve the continuous-variable testing and coding choices, compare methods
given the same admissible knowledge, and keep a numerical-only arm. The
thesis itself acknowledges calibration problems, lost valid edges and the
absence of hidden-confounder handling. Those limitations belong alongside
the favorable reported scores. \citep{rel64molero}

\subsection{Retaining the alternative changes interval coverage}

The two Gaussian mechanisms in Section~\ref{exec20:expert-identification}
make it possible to examine this distinction with a known observation law.
The forward mechanism has intervention effect $3/5$ and the reverse mechanism
has effect zero, while both generate the same standardized joint distribution.
Their ordering is deliberately unresolved; this construction does not reverse
an action and an outcome whose temporal order is known.

A retained project experiment draws 2,000 independent datasets of 500
observations from that common Gaussian law. In each dataset, regression of
the outcome on an intercept and the candidate cause gives an ordinary 95\%
Student interval for the slope. Gaussian independent observations and the
correct linear conditional mean justify that slope interval. It is also an
intervention interval in the forward mechanism. Under the reverse mechanism,
the same calculation estimates a conditional association rather than the
intervention effect.

\begin{table}[htbp]
\centering\small
\begin{tabular}{lrr}
\toprule
Returned set & Forward-world coverage & Reverse-world coverage\\
\midrule
Slope interval & 95.7\% & 0.0\%\\
Slope interval together with zero & 95.7\% & 100.0\%\\
\bottomrule
\end{tabular}
\caption{One observation law and two causal interpretations. Each percentage
uses all 2,000 retained replications, each with 500 observations. The second
set retains both declared mechanisms. This is a generated example, not a
reproduction of neural-flow training or measured banking performance.}
\label{rel69:graph-coverage}
\end{table}

The slope interval covers the forward effect in 1,914 of the 2,000 datasets
and covers the reverse effect in none. Its mean width is 0.1411. The Monte
Carlo standard error of the 95.7\% coverage estimate is 0.45 percentage points;
the observed proportion does not establish improvement over nominal 95\%
coverage. First, middle and last interval calculations were independently
checked by QR decomposition in the original study. Reaggregating all retained
intervals reproduces the counts and width. This checks the stated experiment;
it supplies no calibration result for a fitted neural ensemble.

Adding zero to each slope interval produces a confidence set for the fixed
two-mechanism class. Its forward-world coverage is unchanged, and its
reverse-world coverage is one by construction. The set need not be a single
interval. Taking its convex hull would preserve coverage but obscure the
gap between the two retained causal possibilities.

This is a concrete instance of the graph-retention decomposition in
equation~\eqref{eq:la03-coverage-decomposition}. The true member is retained
with certainty because the example fixes the class to contain both mechanisms.
The forward mechanism contributes its justified regression interval; the
reverse contributes the exact zero effect. For a learned graph class, the
probability of discarding the true member must instead be controlled, and
effect inference must remain valid under the actual selection procedure.
Simultaneous inference over a fixed candidate family or an independent
estimation sample can address particular selection problems. Neither repairs
a false model or supplies missing identification. An unidentified retained
mechanism must contribute justified bounds or an uninformative set; removing
it because it widens the answer introduces an additional error.

The experiment supports a confidence set for its declared two-mechanism
class. General graph-class coverage and bank validation remain open.

\subsection{What the available software can contribute}

The public repository linked by the thesis is a PyWhy-LLM adaptation. Its
2 February 2026 revision contains the language suggestion and retrieval
components, pairwise notebooks, declared dependencies and an MIT notice
carrying the PyWhy copyright. The complete file inventory and two inspected notebooks do not expose the
thesis's HallBayes integration, Monte Carlo graph construction or causal-flow
training pipeline. This is a limit of the available reproduction evidence,
not proof that no private or differently named implementation exists. They support reuse
assessment of the front end, not independent reproduction of all reported
results. Its retrieval helper also creates its own language-model client and
falls back to an answer without retrieved evidence when no match is found.
A bank adapter should pass the selected model explicitly and retain the
absence of evidence as an observable condition. The front-end ideas therefore belong behind a small, versioned interface.
An upstream component that supplies the required operation with maintained
interfaces is a serious alternative to maintaining this adaptation. The
inspected repository establishes no performance advantage for a wholesale
port. Its MIT notice is distinct from the permissions needed for dependencies,
retrieval data and models.
\citep{closure30molerocode}

The HallBayes repository at the inspected June 2026 revision presents
Berry, an evidence-support service with command-line and MCP interfaces.
Its README distinguishes cited-evidence support from world truth, and the
inspected core accepts claims and cited spans as the scoring inputs.
That distinction fits a useful role in checking an assistant's attributed
claims. The June 2026 revision is not the December 2025 implementation cited
by the thesis. A date-bounded query of the current default-branch history
returned no commit for that earlier date, while a retained version-one tag
points to February 2026. This does not prove the older software never
existed; it prevents silently treating the current service as the thesis's
exact dependency. No downloaded source was executed in this comparison.
\citep{closure30berry}

The decision is consequently specific. Adopt the separation between
evidence retrieval, graph construction and intervention calculation.
Evaluate quotation and evidence-support checks for their own stated
predicates. Decline the direct conversion of an information ceiling into
calibrated edge probabilities, and decline unqualified confidence-interval
language for the spread of fitted-model effects. Before a prototype makes
stronger claims, it needs an authentic complete implementation, a declared
joint graph law or retention procedure, the relevant calibration evidence,
and public or synthetic experiments with known interventions and matched
information across comparators. Bank validation can follow when those
experiments justify its cost.

For the bank, the investment conclusion follows from the retained example in
Section~\ref{exec20:expert-identification}. Its observationally identical
mechanisms imply margins of USD~20 and USD~$-40$ under the same stated
valuation and cost. A narrow distribution of repeated fits cannot resolve
that disagreement. Before releasing a financial recommendation, the workflow
must state the decision-time information that excludes a mechanism, or carry
the ambiguity into the recommendation. Graph proposals, retrieval and flexible
conditional models remain candidate components. Their integration deserves
evaluation against an expert-reviewed graph, a numerical discovery route and
an ambiguity-preserving analysis, with withheld mechanisms, interval coverage,
decision loss and complete review cost assessed separately.
